% ch05_u02 —— 原书 p.201–p.209（5.3.3 的后半句 – 5.3.17）
% 首句承接 ch05_u01 末尾的「故按」（原书 p.200 末 "So by"），拼接后读作「故按 Deligne 定理，……」。
\SourcePage{201}

Deligne 定理\footnote{不用 Deligne 定理也可注意到：$\operatorname{Orb}_\chi$ 上任一向量丛的解析截面都是代数的。把它用于带可积联络的向量丛在 $\operatorname{Orb}_\chi$ 上的水平解析截面，即知 (ii) 由 (i) 推出。}，$\operatorname{Orb}_\chi$ 上的光滑 $D$-模具有正则奇点，因此 (ii) 由 (i) 推出。\hfill$\Box$

记 $\mathcal{P}$ 为 $\mathcal{GR}$ 上沿我们的分层光滑的凝聚（等价地，完整）$D$-模构成的范畴。

\bdp{5.3.3} \emph{命题.}

\begin{enumerate}[label=(\roman*),leftmargin=2.2em]
\item 范畴 $\mathcal{P}$ 是半单的。
\item 若 $M\in\mathcal{P}$ 支撑在偶（相应地，奇）连通分支上，则当 $a$ 为奇（相应地，偶）数时 $H^{a}_{DR}(\mathcal{GR},M)=0$。
\end{enumerate}

\emph{证明.} 记 $I_\chi$ 为 $\operatorname{Orb}_\chi$ 上 $\mathbb{C}$-向量空间构成的相交上同调反常层。记 $\mathcal{GR}_{(\chi)}$ 为 $\mathcal{GR}$ 中含 $\operatorname{Orb}_\chi$ 的连通分支，$p(\chi)$ 为 $\mathcal{GR}_{(\chi)}$ 的奇偶性。由 Lusztig（[Lu83] 的定理 11c），$I_\chi$ 具有下述性质：上同调层 $H^{i}(I_\chi)$ 为零，除非 $i\bmod 2=p(\chi)$。记 $C$ 为 $D^{b}(\mathcal{GR}_{(\chi)})$ 中具有该性质、且沿我们的分层光滑的全部对象构成的范畴。由 (259) 与 5.3.2(iii)，对任意 $M,N\in C$，只要 $i\bmod 2\neq p(\chi)$ 就有 $H^{i}(\mathcal{GR}_{(\chi)},M)=0$；且对奇数 $i$ 有 $\operatorname{Ext}^{i}(M,N^{*})=0$（这里 $N^{*}$ 是 $N$ 的 Verdier 对偶）。特别地，只要 $i\bmod 2\neq p(\chi)$ 就有 $H^{i}(\mathcal{GR},I_\chi)=0$；又 $\operatorname{Ext}^{1}(I_{\chi_1},I_{\chi_2})=0$。用 5.3.2(ii) 即得命题。\hfill$\Box$

\bdp{5.3.4} 由 5.3.2(ii)，$\mathcal{P}$ 的单对象是各层 $\operatorname{Orb}_\chi$ 的「相交上同调」$D$-模 $I_\chi$。于是 5.3.3(i) 表明，$\mathcal{P}$ 的任一对象都具有 $G(O)$-等变或 $\operatorname{Aut}^{0}O\ltimes G(O)$-等变的 $D$-模结构。这种结构唯一，且任何态射都与它相容（因为我们的群连通）。我们看到，\SourcePage{202}
$\mathcal{P}$ 等同于 $\mathcal{GR}$ 上 $G(O)$-等变或 $\operatorname{Aut}^{0}O\ltimes G(O)$-等变的凝聚 $D$-模构成的范畴。

\emph{注.} $G(O)$-等变结构的存在性也可直接由下述事实推出：$G(O)$ 连通，且 $\operatorname{Hom}(G(O),\mathbb{G}_m)=0$（以及 5.3.2(ii)）。此时无需援引 5.3.3(i)，因而也无需援引 Lusztig 定理（那是一个深刻的结果）。

\bdp{5.3.5} 范畴 $\mathcal{P}$ 带有典范的张量结构。有两种方式描述它：「卷积」构造（见 5.3.5 -- 5.3.9）与「融合」构造（先作 5.3.10 -- 5.3.12 的若干准备，再在 5.3.13 -- 5.3.16 中给出）；两种定义的等价性见 5.3.17。我们从卷积图景\footnote{以下内容是 Ginzburg 的拓扑构造 [Gi95] 的代数版本；辨认这两种构造的工作留给感兴趣的读者。}开始。我们需要定义卷积积函子 $\circledast:\mathcal{P}\times\mathcal{P}\to\mathcal{P}$、$\circledast$ 的结合约束，以及交换约束。

由 [MV]，函子 $\circledast$ 定义如下。记 $G(K)\times_{G(O)}\mathcal{GR}$ 为 $G(K)\times\mathcal{GR}$ 对 $G(O)$ 的商，其中 $u\in G(O)$ 通过 $(g,x)\mapsto(gu^{-1},ux)$ 作用在 $G(K)\times\mathcal{GR}$ 上。态射 $p:G(K)\times_{G(O)}\mathcal{GR}\to G(K)/G(O)=\mathcal{GR}$，$(g,x)\mapsto g\bmod G(O)$，是纤维为 $\mathcal{GR}$ 的局部平凡纤维化；它附属于主 $G(O)$-丛 $G(K)\to\mathcal{GR}$ 与 $G(O)$ 在 $\mathcal{GR}$ 上的作用。故 $G(K)\times_{G(O)}\mathcal{GR}$ 是 $\mathcal{GR}\times\mathcal{GR}$ 的一个扭形式。设 $M,N\in\mathcal{P}$。用 $M$ 上的 $G(O)$-等变结构，可定义 $G(K)\times_{G(O)}\mathcal{GR}$ 上的一个 $D$-模 $M\boxtimes'N$，它是 $M\boxtimes N$ 的「扭形式」。于是

\begin{equation}\tag{260}
M\circledast N=m_{*}(M\boxtimes'N),
\end{equation}

其中 $m:G(K)\times_{G(O)}\mathcal{GR}\to\mathcal{GR}$ 来自 $G(K)$ 在 $\mathcal{GR}$ 上的作用。

\bdp{5.3.6} \emph{奇迹定理（Lusztig）.} 若 $M,N\in\mathcal{P}$，则 $M\circledast N\in\mathcal{P}$。\hfill$\Box$

\SourcePage{203}

\emph{注.} 非平凡的论断是 $M\circledast N$ 为 $D$-模（而不只是导出范畴的对象）。由于该 $D$-模凝聚且 $G(O)$-等变，它属于 $\mathcal{P}$。

\emph{注.} G. Lusztig 对 5.3.6 的原证明 [Lu83] 基于显式估计，说明 $m$ 限制到 $\mathcal{GR}$ 中两个层的（扭）积上是半小的。很久以后 D. Gaitsgory 指出 [Ga2001]：只要把融合张量积定义中涉及的几何结构考虑进来，5.3.6 便立即成立。

于是我们定义了 $\circledast:\mathcal{P}\times\mathcal{P}\to\mathcal{P}$。$\circledast$ 的结合约束按显然的方式定义。交换约束将在 5.3.8 中定义。

\bdp{5.3.7} \emph{注.} (i) 设 $G(K)$ 换成一个归纳仿射的群归纳概形 $\mathcal{G}$，$G(O)$ 换成其闭子群概形 $\mathcal{K}$；并设 $\mathcal{G}/\mathcal{K}$ 是归纳有限型的归纳概形。5.3.5 中 $\circledast:\mathcal{P}\times\mathcal{P}\to\mathcal{P}$ 的构造基于奇迹 5.3.6。一般地，$\mathcal{G}/\mathcal{K}$ 上 $\mathcal{K}$-等变 $D$-模的范畴上没有卷积；必须考虑某个导出范畴 $\mathcal{H}$（Hecke 幺半范畴；见 7.6.1 与 7.11.17）。这是一个带 t-结构的三角范畴，其核心是 $\mathcal{G}/\mathcal{K}$ 上 $\mathcal{K}$-等变 $D$-模的范畴；一般地，$\circledast:\mathcal{H}\times\mathcal{H}\to\mathcal{H}$ 不是 t-正合的，且 $\circledast$ 没有交换约束。在 $(G(K),G(O))$ 的情形，由 5.3.6，函子 $\circledast$ 是 t-正合的，而 $\mathcal{H}$ 的核心是 $\mathcal{P}$ 的归纳对象构成的范畴。

(ii) 上述 $\mathcal{H}$ 的构造是 §7 中发展的「Hecke 模式」的一部分。后面我们将看到，即使在 $(G(K),G(O))$ 这一奇迹般好的情形，该模式也是有用的（或许不可或缺）。

\bdp{5.3.8} 我们来定义 $\circledast$ 的交换约束。设 $\theta:G\to G$ 是把每个支配权映到其对偶的自同构。$G$ 的反自同构 $\theta'(g):=\theta(g)^{-1}$ 给出 $\mathcal{H}$ 的反自同构 $\theta'_{\mathcal{H}}$\SourcePage{204}，故对任意 $M,N\in\mathcal{H}$ 有典范同构

\[
l_{M,N}:\theta'_{\mathcal{H}}(M\circledast N)\xrightarrow{\sim}\theta'_{\mathcal{H}}(N)\circledast\theta'_{\mathcal{H}}(M).
\]

对任意 $M\in\mathcal{P}\subset\mathcal{H}$ 有典范同构 $e_M:M\xrightarrow{\sim}\theta'_{\mathcal{H}}(M)$。按 5.3.3(i)，为定义 $e_M$ 只需考虑 $M=I_\chi$ 的情形。$\theta'$ 在 $G(K)$ 上的作用保持由双 $G(O)$-类给出的分层 $G(K)_\chi$（这里 $G(K)_\chi$ 是 $\operatorname{Orb}_\chi\subset G(K)/G(O)$ 的原像）。于是有诱导的自同构 $\theta'_\chi$。作为 $\mathcal{H}$ 的对象，我们的 $I_\chi$ 是 $G(K)$ 上的 $\Omega$-复形 $\Omega_{G(K)_\chi}[\dim\operatorname{Orb}_\chi]$。现在 $e_{I_\chi}$ 是 $\theta'_\chi$ 在 $\Omega_{G(K)_\chi}$ 上的作用。

对 $M,N\in\mathcal{P}$ 定义

\begin{equation}\tag{261}
s:M\circledast N\xrightarrow{\sim}N\circledast M
\end{equation}

为如下复合：

\[
M\circledast N\xrightarrow{\sim}\theta'_{\mathcal{H}}(M\circledast N)\xrightarrow{\sim}\theta'_{\mathcal{H}}(N)\circledast\theta'_{\mathcal{H}}(M)\xrightarrow{\sim}N\circledast M,
\]

其中第一个箭头是对应于 $M\circledast N$ 的同构 $e$，其余箭头是 $l_{M,N}$ 与 $e_N^{-1}\circledast e_M^{-1}$。

\bdp{5.3.9} \emph{命题.} $s$ 是卷积张量积 $\circledast$ 的交换约束。

\emph{证明.} 在下文 5.3.17 中，我们把卷积张量积与融合张量积等同起来，且该等同与所有约束相容。由于后一组数据显然在 $\mathcal{P}$ 上定义了张量范畴结构，证毕。\hfill$\Box$

于是我们定义了 $\mathcal{P}$ 上预期的卷积张量结构。

\bdp{5.3.10} $\mathcal{P}$ 上张量结构的融合描述\footnote{该构造表面上涉及一条曲线 $X$，但实际上纯粹是局部的。}基于 4.3.14 中「空间」$\operatorname{GRAS}=\operatorname{GRAS}_G$ 上重要的手性半群结构。该结构可描述如下。

\SourcePage{205}

(i) 对 $\mathbb{C}$-代数 $R$ 与 $S\in\Sigma(R)$（我们沿用 4.3.11 的记号，故 $S$ 是 $X\otimes R$ 的一个在 $\operatorname{Spec}R$ 上有限且平坦的子概形），有子集 $\operatorname{GRAS}(R)_S\subset\operatorname{GRAS}(R)$，它由对 $(\mathcal{F},\gamma)$ 构成，其中 $\mathcal{F}$ 是 $X\otimes R$ 上的 $G$-丛，$\gamma$ 是 $\mathcal{F}$ 在 $S$ 的补集上的截面。

(ii) 若 $S$ 是子概形 $S_i$（$i\in I$）的不交并，则有典范等同

\begin{equation}\tag{262}
\operatorname{GRAS}(R)_S\xrightarrow{\sim}\prod_i\operatorname{GRAS}(R)_{S_i}.
\end{equation}

确切地说，我们把 $(\mathcal{F},\gamma)$ 与族 $(\mathcal{F}_i,\gamma_i)$（$i\in I$）等同，其中 $(\mathcal{F}_i,\gamma_i)\in\operatorname{GRAS}(R)_{S_i}$ 在 $S_{i'}$（$i'\neq i$）之并的补集上与 $(\mathcal{F},\gamma)$ 一致。

数据 (i)、(ii) 具有如下性质：

a. 若对 $S_1,S_2\in\Sigma(R)$ 有 $S_{1\mathrm{red}}\subset S_{2\mathrm{red}}$，则 $\operatorname{GRAS}(R)_{S_1}\subset\operatorname{GRAS}(R)_{S_2}$。$\operatorname{GRAS}(R)_S$（$S\in\Sigma(R)$）之并等于 $\operatorname{GRAS}(R)$。故 $\operatorname{GRAS}(R)_S$ 构成 $\operatorname{GRAS}(R)$ 上的滤链。该滤链具有函子性（关于 $R$）。

b. 同构 (ii) 也具有函子性，并按显然的方式与 $I$ 的细分相容。

c. 由

\[
\mathcal{GR}_\Sigma(R):=\{(S,\mathcal{F},\gamma)\mid S\in\Sigma(R),\ (\mathcal{F},\gamma)\in\operatorname{GRAS}(R)_S\}
\]

定义的子函子 $\mathcal{GR}_\Sigma\subset\Sigma\times\operatorname{GRAS}$ 是 $\Sigma$ 上形式光滑的归纳概形。

\emph{注.} 我们来解释对任意仿射代数群 $G$，$\mathcal{GR}_\Sigma=\mathcal{GR}^{G}_\Sigma$ 为何是归纳概形。此外我们将证明：$\mathcal{GR}_\Sigma$ 是归纳有限型的；且当 $G$ 约化时 $\mathcal{GR}_\Sigma$ 归纳固有。先考虑 $G=GL_n$ 的情形。此时 $\mathcal{GR}_\Sigma$ 是诸 $\mathcal{GR}_{\Sigma,k}$ 的正向极限，其中 $\mathcal{GR}_{\Sigma,k}$ 参数化由有限子概形 $D\subset X$ 与子层 $\mathcal{E}\subset\mathcal{O}_X^{n}(kD)$ 构成、且满足 $\mathcal{E}\supset\mathcal{O}_X^{n}(-kD)$ 的对。态射 $\mathcal{GR}_{\Sigma,k}\to\Sigma$ 固有，故 $\mathcal{GR}_\Sigma$ 归纳固有。如定理 4.5.1 的证明中所说明，为把一般情形化到\SourcePage{206} $GL_n$ 的情形，只需证明：若 $G\subset G'$ 且 $G'/G$ 仿射（相应地，拟仿射），则态射 $\mathcal{GR}^{G}_\Sigma\to\mathcal{GR}^{G'}_\Sigma$ 是闭（相应地，局部闭）嵌入。这很容易。

\bdp{5.3.11} 对有限集 $J$，有态射 $X^{J}\to\Sigma$，它把 $(x_j)\in X^{J}$ 对应到与除子 $\sum_j x_j$ 相应的子概形 $D\subset X$。记 $\mathcal{GR}_{X^{J}}$ 为 $\mathcal{GR}_\Sigma$ 与 $X^{J}$ 在 $\Sigma$ 上的纤维积。于是 $\mathcal{GR}_{X^{J}}$ 的一个 $R$-点是一个族 $((x_j),\mathcal{F},\gamma)$，其中 $(x_j)\in X^{J}(R)$，$\mathcal{F}$ 是 $X\otimes R$ 上的 $G$-丛，$\gamma$ 是 $\mathcal{F}$ 在诸 $x_j$ 的图像之并的补集上的截面。我们的 $\mathcal{GR}_{X^{J}}$ 是 $X^{J}$ 上形式光滑、归纳固有的归纳概形（见 5.3.10 末的注）。

由 4.5.2，$\mathcal{GR}_X$ 在 $x\in X(\mathbb{C})$ 上的纤维与归纳概形 $\mathcal{GR}_x:=G(K_x)/G(O_x)$ 之间有典范同构。故由 5.3.10(ii)，$\mathcal{GR}_{X^{J}}$ 在 $(x_j)\in X^{J}(\mathbb{C})$ 上的纤维等于 $\prod_{x\in S}\mathcal{GR}_x$，其中 $S$ 是子集 $\{x_j\}\subset X$。

下面的 $\mathcal{GR}_X$ 的描述将有用。考虑 $X$ 上「形式参数」的概形 $X^{\wedge}$（其点是光滑态射 $\operatorname{Spec}O\to X$，见 2.6.5）。这是 $X$ 上的 $\operatorname{Aut}^{0}O$-torsor；坐标的选择，即某个开集 $U\subset X$ 上平展的 $\mathbb{A}^{1}$-值映射，给出 $X^{\wedge}$ 在 $U$ 上的平凡化。现在 $\mathcal{GR}_X$ 是 $\mathcal{GR}$ 的 $X^{\wedge}$-扭（关于 $\operatorname{Aut}^{0}O$ 在 $\mathcal{GR}$ 上的作用）。$\mathcal{GR}$ 的分层给出 $\mathcal{GR}_X$ 的一个分层，其层为 $X$ 上光滑的 $\operatorname{Orb}_{\chi_X}$。

\bdp{5.3.12} 为便于将来引用，我们列出 $\mathcal{GR}_{X^{J}}$ 之间的一些相容性；它们直接由 5.3.10 推出。

a. 对满射 $\pi:J\twoheadrightarrow J'$，有显然的笛卡尔图表

\begin{equation}\tag{263}
\begin{tikzcd}
\mathcal{GR}_{X^{J'}} \arrow[r,"\tilde\Delta^{(\pi)}"] \arrow[d] & \mathcal{GR}_{X^{J}} \arrow[d]\\
X^{J'} \arrow[r,"\Delta^{(\pi)}"] & X^{J}
\end{tikzcd}
\end{equation}

\SourcePage{207}

其中 $\Delta^{(\pi)}$ 是 $\pi$-对角嵌入。若 $|J'|=1$，则有 $\Delta^{(J)}:X\hookrightarrow X^{J}$ 与 $\tilde\Delta^{(J)}:\mathcal{GR}_X\hookrightarrow\mathcal{GR}_{X^{J}}$。

b. 设 $\nu^{(J)}:U^{(J)}\hookrightarrow X^{J}$ 是对角除子的补集。由 5.3.10(ii)，$X^{J}$-归纳概形 $\mathcal{GR}_{X^{J}}$ 与 $(\mathcal{GR}_X)^{J}$ 在 $U^{(J)}$ 上的限制有典范等同。因此有笛卡尔图表

\begin{equation}\tag{264}
\begin{tikzcd}
(\mathcal{GR}_X)^{J}|_{U^{(J)}} \arrow[r,"\tilde\nu^{(J)}"] \arrow[d] & \mathcal{GR}_{X^{J}} \arrow[d]\\
U^{(J)} \arrow[r,"\nu^{(J)}"] & X^{J}
\end{tikzcd}
\end{equation}

\bdp{5.3.13} 现在我们准备定义 $\mathcal{P}$ 上的融合张量结构。这归结为对任意有限非空集 $J$ 构造张量积函子\footnote{这里 $\mathcal{P}^{\otimes J}$ 表示 $\mathcal{P}$ 的 $J$ 个副本的张量积（由于 $\mathcal{P}$ 半单，张量积的定义是清楚的）。}

\begin{equation}\tag{265}
\circledast_J:\mathcal{P}^{\otimes J}\to\mathcal{P}
\end{equation}

并给出等同

\begin{equation}\tag{266}
\circledast_J=\circledast_{J'}\Bigl(\bigotimes_{j'\in J'}\bigl(\circledast_{\pi^{-1}(j')}\bigr)\Bigr)
\end{equation}

对任意满射 $J\xrightarrow{\pi}J'$。

构造如下。

\bdp{5.3.14} 由于任一 $M\in\mathcal{P}$ 是 $\operatorname{Aut}^{0}O$-等变的，它定义了 $\mathcal{GR}_X$ 上的一个 $D$-模（见 5.3.11 末对 $\mathcal{GR}_X$ 的描述）。记 $M_X\in D(\mathcal{GR}_X)$（$:=DM(\mathcal{GR}_X)$）为它在导出范畴中的平移 1。换言之，对上述任一开集 $U$ 与 $X^{\wedge}$ 在 $U$ 上的平凡化 $\theta$，有 $M_U=\pi_\theta^{!}M$；其中 $M_U:=M_X|_{\mathcal{GR}_U}$，$\pi_\theta:\mathcal{GR}_U\to\mathcal{GR}$ 是对应于 $\theta$ 的投影。我们用 $M$ 上的 $\operatorname{Aut}^{0}O$-作用把这些对象粘合起来。函子 $\mathcal{P}\to D(\mathcal{GR}_X)$，$M\mapsto M_X$，是全忠实的。其本质像由（平移 1 的）$D$-模构成，这些 $D$-模同构于「相交上同调」$D$-模 $I_{\chi X}$ 的（有限多个）副本的直和；$I_{\chi X}$ 对应于 $\operatorname{Orb}_{\chi X}$ 上的平凡局部系统。

\SourcePage{208}

设 $\{M_j\}_{j\in J}$ 是 $\mathcal{P}$ 中对象构成的族。用 (264)，把 $\boxtimes M_{jX}|_{U^{(J)}}$ 解释为 $\mathcal{GR}_{X^{J}}|_{U^{(J)}}$ 上的、平移 $|J|$ 的 $D$-模。记 $\boxtimes M_{jX}\in D(\mathcal{GR}_{X^{J}})$ 为它到 $\mathcal{GR}_{X^{J}}$ 的极小（即 $\tilde\nu^{(J)}_{!*}$-）扩张。这是 $\mathcal{GR}_{X^{J}}$ 上平移 $|J|$ 的 $D$-模。因此我们定义了函子

\begin{equation}\tag{267}
\boxtimes_J:\mathcal{P}^{\otimes J}\to D(\mathcal{GR}_{X^{J}}),\qquad \otimes M_j\mapsto\boxtimes M_{jX},
\end{equation}

它显然是全忠实的。

\bdp{5.3.15} \emph{命题.}（[MV]）

对任意 $\pi:J\twoheadrightarrow J'$，复形 $\tilde\Delta^{(\pi)!}(\boxtimes M_{jX})\in D(\mathcal{GR}_X^{(J')})$ 属于 $\boxtimes_{J'}$ 的本质像。\hfill$\Box$
% [Tutor Note: 原书 p.208 此处印作 $D(\mathcal{GR}_X^{(J')})$，与 (267) 中的 $\mathcal{GR}_{X^{J}}$ 写法不一致；底本如此，保留原样。]

\bdp{5.3.16} 我们得到函子

\begin{equation}\tag{268}
\circledast_\pi:\mathcal{P}^{\otimes J}\to\mathcal{P}^{\otimes J'},
\end{equation}

使得 $\boxtimes_{J'}\circledast_\pi=\tilde\Delta^{(\pi)!}\boxtimes_J$。特别地，当 $|J'|=1$ 时得到函子 $\circledast_J:\mathcal{P}^{\otimes J}\to\mathcal{P}$，它就是我们的张量积函子 (265)。显然的等同 $\circledast_\pi=\bigotimes_{j'\in J'}\bigl(\circledast_{\pi^{-1}(j')}\bigr)$（看我们在 $U^{(J')}$ 上的 $D$-模）与标准同构 $\Delta^{(J)!}=(\Delta^{(\pi)}\Delta^{(J')})^{!}=\Delta^{(J')!}\Delta^{(\pi)!}$ 给出相容性同构 (266)。故 $\mathcal{P}$ 是张量范畴。容易看出 $I_0$ 是 $\mathcal{P}$ 中的单位对象。

\bdp{5.3.17} 我们来把 $\mathcal{P}$ 上的卷积与融合张量结构等同起来。下面在本小节中，用 $\circledast^{c}$ 表示卷积张量积，用 $\circledast^{f}$ 表示 $\mathcal{P}$ 上的融合张量积。我们要对 $M,N\in\mathcal{P}$ 构造与结合约束、交换约束相容的典范同构 $M\circledast^{c}N\xrightarrow{\sim}M\circledast^{f}N$。\footnote{该构造借自 [MV]，那里写得更详细；不过那里没有考虑交换约束 5.3.8。}

设 $\mathcal{GR}'_{X^{2}}$ 是 $X^{2}$ 上的归纳概形，使得 $\mathcal{GR}'_{X^{2}}(R)$ 是族 $(x_1,x_2,\mathcal{F}_1,\mathcal{F}_2,\gamma_1,\gamma_2)$ 的集合，其中 $x_1,x_2\in X(R)$，$\mathcal{F}_1,\mathcal{F}_2$ 是 $X\otimes R$ 上的 $G$-torsor，$\gamma_1$ 是 $\mathcal{F}_1$ 在\SourcePage{209} $x_1$ 的图像之补集上的截面，$\gamma_2$ 是 $\mathcal{F}_1\to\mathcal{F}_2$ 在 $x_2$ 的图像之补集上的同构。有投影 $q:\mathcal{GR}'_{X^{2}}\to\mathcal{GR}_{X^{2}}$，它把上述数据映到 $((x_1,x_2),\mathcal{F}_2,\gamma_2\gamma_1)$。该投影是归纳固有的；在 $U:=X^{2}\setminus\{\text{对角线}\}$ 上它是同构。\footnote{在对角线上，$q$ 的纤维同构于 $\mathcal{GR}$；更确切地说，由 $(x_1,x_2,\mathcal{F}_1,\mathcal{F}_2,\gamma_1,\gamma_2)\mapsto(x_1,x_2,\mathcal{F}_1,\gamma_1,\mathcal{F}_2,\gamma_2\gamma_1)$ 定义的闭嵌入 $\mathcal{GR}'_{X^{2}}\to(\mathcal{GR}_X)\times_X(\mathcal{GR}_{X^{2}})$ 限制到对角线 $X\hookrightarrow X^{2}$ 上成为同构。故 $q$ 为同构的最大开子集形如 $\mathcal{GR}_{X^{2}}\setminus Z$，其中 $Z$ 的余维为 1；这是一个无穷维现象。}
% [Tutor Note: 原书 p.209 脚注中闭嵌入的目标印作 $(\mathcal{GR}_X)\times_X(\mathcal{GR}_{X^2})$，按上下文疑为 $(\mathcal{GR}_X)\times_X(\mathcal{GR}_X)$ 之误；底本如此，保留原样。]

记 $M_X\boxtimes'N_X\in D(\mathcal{GR}'_{X^{2}})$ 为 $M_X\boxtimes N_X|_U$ 到 $\mathcal{GR}'_{X^{2}}$ 的极小扩张。这是 $\mathcal{GR}'_{X^{2}}$ 上平移 2 的 $D$-模。由 [MV]，$U$ 上的显然等同（唯一地）延拓为典范同构

\begin{equation}\tag{269}
q_{*}(M_X\boxtimes'N_X)\xrightarrow{\sim}M_X\boxtimes N_X.
\end{equation}

现在 $\mathcal{GR}'_{X^{2}}$ 是 $(\mathcal{GR}_X)^{2}$ 的扭形式。事实上，$\mathcal{F}_1$ 在 $x_2$ 的形式邻域上的平凡化给出上述数据 $(\mathcal{F}_2,\gamma_2)$ 与 $\mathcal{GR}_x$ 的一个等同。这些平凡化连同 $x_2$ 处的形式参数构成 $\mathcal{GR}_X\times X$ 上的 $\operatorname{Aut}^{0}O\ltimes G(O)$-torsor，而 $\mathcal{GR}'_{X^{2}}$ 与 $\mathcal{GR}$ 的相应扭等同。故 $M_X\boxtimes'N_X$ 是 $M_X\boxtimes N$ 的「扭形式」。把这一图景限制到对角线 $X\hookrightarrow X\times X$，我们看到：$q:\mathcal{GR}'_{X^{2}}\to\mathcal{GR}_{X^{2}}$ 到 $X$ 的拉回与 (260) 中态射 $m:G(K)\times_{G(O)}\mathcal{GR}\to\mathcal{GR}$ 的 $X^{\wedge}$-扭重合；而 $M_X\boxtimes'N_X$ 到 $X$ 在 $\mathcal{GR}'_{X^{2}}$ 中前像的拉回等于 $(M\boxtimes'N)_X$，其中 $M\boxtimes'N$ 的含义同 (260)。比较 (269) 与 (260)（并用基变换同构），即得所要的典范同构 $M\circledast^{c}N\xrightarrow{\sim}M\circledast^{f}N$。

它与结合约束的相容性来自 $X^{3}$ 上的类似图景。写出与交换约束的相容性（用 $\operatorname{Bun}_G$ 与 Hecke）！
% [Tutor Note: 原书 p.209 末句为作者的大写备忘（"WRITE DOWN THE COMAT WITH COM CONSTRAINTS"，即 compatibility with commutativity constraints）；底本如此，保留原样。]
